\cotefile{157-5}
\scanpage{528}
\mspage{439}
% archive sheet number: 527

\deleted{\unsure{Passons maintenant que|Prouvons maintenant que|Posons maintenant que}}

Dans le cas où $f$ est fibrant, admettons (comme hyp. suppl au besoin) que $\deleted{\unsure{f_s|f}},\unsure{E|\mathcal E}$ satisfasse $L$ et que $f$ soit accessible, et prenons $\pi$ assez grd pour que $\deleted{\unsure{f|E_s}}E,f$ satisfassent $L_\pi$. Nous pouvons alors terminer la démonstration de 4.15.6. Il faut prouver

\[ E(\pi)\subset E_\pi \]

donc que si $x\in\mathrm{Ob}\,E_{s\pi}$, et \unsure{considérant|considérons} un système inductif dans $E$ grd devant $\pi$, de \uline{lim} $y$, on a

\[ \varinjlim\mathrm{Hom}_{E}(x,y_i)\xrightarrow{\sim}\mathrm{Hom}_{E}(x,y). \]

\deleted{\unsure{O|D}}Donc il faut prouver que tte flèche $x\xrightarrow{u}y$ se factorise, de façon ess. unique\margin{\unsure{unq|uniq}}, par un $y_i$. Mais la flèche $s\xrightarrow{\lambda=f(u)}t$ dans $C$ se factorise, de façon ess. unique, par des $\lambda_i:s\mathrel{\unsure{\relbar|\to}}t_i$. Donc on est ramené à trouver \deleted{\unsure{un|une}} \uline{$\lambda_i$-morphisme} $x\longrightarrow y_i$ pour $i$ assez grand qui factorise. Dans le cas fibrant, cela signifie qu'on cherche $x\longrightarrow\lambda_i^*(y_i)$ qui factorise $x\longrightarrow\lambda^*(y)$, or les $\lambda_i^*(y_i)$ forment un système inductif \deleted{dans} $E_s$, et sa limite \uline{dans $E_s$} est $\deleted{\unsure{s|t}}\lambda^*(y)$ (\uline{à vérifier}). Comme $x$ est de $E_{s\pi}$, l'hom donné $x\longrightarrow\lambda^*(y)$ se factorise de
\scanpage{529}%
\mspage{440}%
% archive sheet number: 528
de façon \unsure{ess\textsuperscript{t}.|esst.} unique par un $\lambda_i^*(y_i)$.

Cette fois par contre, c'est le cas cofibrant qui pose problème. En fait, je n'arrive pas à m'en tirer.

Ainsi, on trouve bien un foncteur

\[ \mathrm{Ind}_\pi\bigl(\unsure{E|\mathcal E}(\pi)\bigr)\longrightarrow E, \]

qui est ess. surjectif, mais faute de savoir que $E(\pi)\subset E_\pi$ dans le cas cofibrant, il n'est pas clair que ce foncteur soit une équivalence de catégories.

4.15.7. Il me faut \unsure{reporter|reprendre} la démonstration du lemme 4.15.3\inserted{(p. \unsure{428|420})}, en admettant 4.8.13. (\unsure{Non encore démontré|Non encore démontrée} -- Je vais revenir plus bas sur la dém. de 4.8.13), On va utiliser le

\uline{Lemme} 4.15.7.1. Soi\deleted{\unsure{ent|nt}}t $C$ une catégorie, $\underline{x}=(x_i)_{i\in I}$ un Ind-objet de $C$. \deleted{\unsure{a)|g}}\deleted{qui soit}\inserted{Pour} que \deleted{$\unsure{C^\natural|C}$ une enveloppe de Karoubi de $C$, $x_0$ un objet de $C^{\deleted{\unsure{\natural|E}}}$ tel que \deleted{\unsure{qui|que}} $\underline{x}$ soit isomorphe dans $\mathrm{Ind}(C)$,} $\underline{x}$ soit \inserted{ess. constant, i.e.} isomorphe, dans $\mathrm{Ind}(C)$, à un \deleted{\unsure{un|m}} Ind-objet constant \deleted{\unsure{et|d}} défini par un objet \deleted{\illegible{1}}
\scanpage{530}%
\mspage{441}%
% archive sheet number: 529
de $C$, il f. et s. qu'il existe un foncteur cofinal $\varphi:\deleted{\unsure{I}}J\to I$, avec $\deleted{I}J$ filtrant, de telle façon que le système réindexé $(x_{\deleted{\unsure{i}}j})^{\deleted{\unsure{i}}}_{j\in J}$, avec $x_j=x_{\varphi(j)}$, ait la propriété suivante : pour tt $j\in J$, on peut trouver un projecteur $p_j$ dans $x_j$, de telle façon que : $1^\circ)$ Pour tte flèche $\lambda:j\mathrel{\relbar}j'$ dans $J$, on ait commutativité dans

\diagram{\begin{tikzcd}[ampersand replacement=\&]
x_j \arrow[r,dash,"{\deleted{\unsure{\tau_{j'j}}}\,\inserted{x_\lambda}}"] \arrow[d,"{p_j}"'] \& x_{j'} \arrow[d,"{p_{j'}}"] \\
x_j \arrow[r,dash,"{\deleted{\unsure{\tau_{j'j}}}\,\inserted{x_\lambda}}"'] \& x_{j'}
\end{tikzcd}}

i.e. $p_{j'}x_\lambda \unsure{g|j}=\deleted{\unsure{\tau_{j'j}}}\inserted{x_\lambda}p_j$ (où les $\tau_{j',j}$ sont les flèches de transition) ; \deleted{\unsure{2)}}

$2^\circ)$ pour \deleted{tout} $j\in\mathrm{Ob}\,J$, $\exists$ flèche $\lambda:j\longrightarrow j'$ dans $J$ telle que $p_{j'}x_\lambda=x_\lambda p_j=x_\lambda$ ; $3^\circ)$ Pour \deleted{\unsure{tout}} $j\in\mathrm{Ob}\,J$, le projecteur $p_j$ de $x_j$ admet une image $x_j^{\unsure{\natural|\sharp}}$ ; $4^\circ)$ \deleted{\unsure{telle|l'hom}} pour tte flèche $\lambda\inserted{:j\to j'}$ dans $I$, la flèche $\deleted{\unsure{x_\lambda:x_j|ax:x_j}}x_\lambda^\natural:x_j^\natural\to x_{j'}^\natural$, induite par $x_\lambda$ (grâce à $1^\circ)$ est un isomorphisme.

De plus \deleted{\unsure{si}} on peut supposer $J$ \deleted{\unsure{filtrante}}\inserted{de la} forme $i_0\backslash I$, le foncteur cofinal $J\to I$

% slanted
\margin{étant le foncteur canonique $i_0\backslash I\to I$ ($i_0$ étant un objet fixé, arbitraire de $I$)}

\deleted{b) Supposons que $I$ soit grande devant $\pi$, $\pi$ est un cardinal infini $\geq\operatorname{card}\operatorname{\unsure{red|réel|nerf}} C$. Alors on peut choisir $J$ ordonné, grand devant $\pi$, et telle que l'on ait de plus la condition $5^\circ)$ Tous les \deleted{$p_j$ sont égaux à un} \inserted{$x_j$ sont égaux à} \inserted{un} même \deleted{projecteur \unsure{$p$ de|$p$ et}} objet $x_0$, tous les $p_j$ sont égaux à un même projecteur $p$, et pour}
\scanpage{531}
\mspage{442}
% archive sheet number: 530

Notons tout de suite deux corollaires

\uline{Corollaire} 4.15.7.2. Supposons que $C$ soit sans projecteurs. Alors $\underline{x}=(x_i)_{i\in \unsure{J|I}}$ est un ind-objet ess. constant sss il existe un foncteur cofinal $J\longrightarrow I$, avec $J$ \deleted{\unsure{ordonné}} \inserted{filtrant}, de telle façon que le système inductif réindexé $(x_j)_{j\in J}$ soit \deleted{isomorphe (\uline{dans la catégorie} $\underline{\mathrm{Hom}}(J,C)$) à un système inductif} \uline{constant} (\unsure{défini} par un foncteur constant $J\longrightarrow C$ ; tous les $x_j$ égaux à un $=x$, et les morphismes de \deleted{\unsure{trans|tous}} tous des identités).
% A broken line runs under the crossed wording from isomorphe to inductif.

% slanted
\margin{\unsure{et OPS|cf. OPS|cf. OP 3} $J=i_0\backslash I\xrightarrow[\mathrm{can.}]{} I,$}

La suffisance est claire, pour la nécessité on applique 4.15.7.1, qui nous dit qu'on peut réindexer de façon à satisfaire aux conditions $1^\circ)$ à $4^\circ)$ (\deleted{pour} des projecteurs $p_j$ convenables \unsure{des|de les} $x_j$). Mais l'hyp. sur $C$ implique que \deleted{les} $p_j \unsure{=|\Rightarrow}\mathrm{id}_{x_j}\quad\forall j\in\mathrm{Ob}\,J$, et les conditions $1^\circ,2^\circ,3^\circ$ sont satisfaites tautologiquement pour un tel choix, la condition $4^\circ$ signifie que les morphismes de transition sont des isomorphismes.

\deleted{\unsure{OPS|OP 3} que $J$ soit ordonné. Mais il est immédiat \deleted{pour} -- cela et pour un système inductif $(x_j)_{j\in J}$ donné, avec $J$ \uline{1-connexe}, \inserted{(\unsure{p. ex.|cf. car} $J$ \unsure{filtrant|filtr.})}, la condition que les $x_\lambda$ soient des iso est néc. et suff. pour que ce syst. soit isom. à un syst. ind. constant, d'où la conclusion}

% slanted
\margin{\unsure{Garder les traits biffés|Garder le texte biffé}}
\scanpage{532}
\mspage{443}
% archive sheet number: 531

\uline{Corollaire} 4.15.7.3. Supposons que $I$ soit grd devant $\pi_0$, où $\pi_0$ est un cardinal infini tel que $\pi_0\geq\operatorname{card}\deleted{\unsure{\mathrm{red}|\mathrm{nerf}|\mathrm{morph.}}} C$. Alors $(x_i)_{i\in I}$ est ess.t constant sss on peut trouver un foncteur cofinal $J\longrightarrow I$, avec $J$ \deleted{\unsure{ordonné et}} grand devant $\pi_0$\inserted{(et \unsure{réduit si on veut|réordonné})}, de telle façon que le système réindexé $(x_j)_{j\in J}$ satisfasse les conditions suivantes, \deleted{il existe $1)$ $\forall j\in\mathrm{Ob}\,J$}

% slanted
\margin{La démonstration \unsure{légèrement réarrangée|convenablement réarrangée} montrera qu'on peut prendre pour $J$ une sous-catégorie (pas pleine) d'une des catégories $i_{\unsure{1|0}}\backslash I$, pour un objet convenable $i_1$ \deleted{\unsure{$j$}} de $I$. \deleted{\illegible{1}} On a donc $\operatorname{card}J\leq\pi$ si $\operatorname{card}I\leq\pi$, $\pi$ étant un cardinal infini donné. \deleted{\unsure{Lorsque $I$ admet} \illegible{2}}}

\deleted{$x_j=\unsure{x_0}$} $1^\circ)$ L'application $j\mapsto x_j$ de $\mathrm{Ob}\,J$ dans $\mathrm{Ob}\,C$ est constante, i.e. il existe un $x_0\in\mathrm{Ob}\,C$ tel que $x_j=x_0\quad\forall j\in\mathrm{Ob}\,J$ ; $2^\circ)$ il existe un projecteur $p$ dans $x_0$ tel que pour $j<j'$ dans $J$, on ait $\tau_{j',j}=p$ ; $3^\circ)$ Le projecteur $p$ admet une image dans $C$.

\uline{Dém} : Ces conditions impliquent \inserted{évidemment} celles de 4.15.7.1, en prenant $p_j=p\quad\forall j\in\mathrm{Ob}\,J$, du moins si $J$ n'admet pas de plus grd élément (ce qui fait une difficulté pour $2^\circ)$). Mais si $J$ admet un objet final, \inserted{\unsure{il en est de même de $I$, et|il en est de même de $I$, it}} il est clair \inserted{alors} que \deleted{\unsure{tt}} syst. inductif indexé par $\deleted{J}I$ \deleted{\unsure{admet}} est ess. constant. Donc, dans \deleted{les deux} tous les cas, les conditions $1^\circ)$ $2^\circ)$ $3^\circ)$ ci-dessus impliquent que $\underline{x}$ est ess. constant. Inverse%
\scanpage{533}%
\mspage{444}%
% archive sheet number: 532
ment, supposons que $\underline{x}$ soit ess. constant, \deleted{\unsure{$\exists$|$\forall$}} \deleted{\unsure{$A$}} et \unsure{prouvons|prenons} \inserted{\unsure{que|pour}} \deleted{\unsure{le cri|la cat.}} \unsure{l'on|l'on y} peut réindexer comme indiqué$^{*}$.

% margin note written sideways
\margin{Le cas où $I$ admet un objet final est trivial, \unsure{on prendra|prendre} $J=e$, $p=\mathrm{id}_{x_i}$ Supposons donc \unsure{que $I$ n'ait pas d'objet final|qu'on n'ait pas d'objet final}.}

Comme $I$ grd devant $\pi_0$ et $\pi_0\geq\operatorname{card}\operatorname{Ob} C$, on voit qu'il existe un \inserted{$x_0\in\operatorname{Ob} C$ tel que la} sous-catégorie pleine $I'\inserted{=I_{x_0}}$ de $I$, \deleted{cofinale} formée des $i$ tels que $x_i=x_0$, soit cofinale ds $I$. \deleted{Donc} Comme $(x_i)_{i\in I}$ est ess. constant, il en est de $=$ de $(x_{i'},i'\in I')$, et on peut lui appliquer le critère 4.15.7.1, avec $J\to I'$ cofinal. La démonstration montre de plus que, si $I'$ est grd devant $\pi_0$ (\inserted{cardinal infini} sans autres conditions) OPS $J$ grde devant $\pi_0$. \inserted{De plus, il est clair qu'on peut supposer $J$ ordonné.}

% slanted
\margin{\unsure{en tenir compte dans la dém. de|en tenir compte ds la dém. de} 4.15.7.1.}

Comme $\pi_0\geq\operatorname{card}\operatorname{Hom}(x_0,x_0)$, on voit aussi qu'il existe \deleted{\unsure{un $p\in$|un $p$}} \inserted{un projecteur} $p$ ds $x_0$ tel que la partie \deleted{cofinale} $\inserted{J_p=}J'$ de $J$, \deleted{telle que} formée des $j$ tels que $p_j=p$, soit cofinale. Étant cofinale, elle est égalemt grde devant $\pi_0$. Donc OPS que tous les $p_j$ sont égaux à un même $p$. Les conditions $1^\circ$--$4^\circ$ deviennent

\[1^\circ)\quad p\,\tau_{j'j}=\tau_{j'j}\,p\quad\text{pour }j\leq j'.\]

\[2^\circ)\quad\deleted{p\tau_{j'j}=\tau_{j'j}}\quad\forall j,\ \text{existe }j'\geq j\ \text{avec}\quad p\tau_{j'j}=\tau_{j'j}p=\tau_{j'j}.\]

$3^\circ)$ \deleted{\unsure{$\forall$}} L'image \inserted{$x_0^\natural$} du projecteur $p$ dans $x_0$ existe.

$4^\circ)$ L'endom. $\tau_{j'j}^\natural$ de $x_0^\natural$ est iso, pour $j\leq j'$.
\scanpage{534}
\mspage{445}
% archive sheet number: 533

Or on a le

\uline{Lemme} 4.15.7.4. Soient $C$ une catégorie, $x_0$ un objet de $C$, $p$ un projecteur de $x_0$, $\tau$ un endomorphisme de $x_0$ satisfaisant les relations

\begingroup\makeatletter\tagsleft@true\makeatother
\[\tau p=p\tau=\tau,\tag*{$\unsure{(a)|(*)}$}\]
\endgroup

et tel que l'endomorphisme de $x_0^\natural\overset{\mathrm{déf}}{=}\operatorname{Im}p$ (défini p. ex. comme objet de $C^\wedge$, $x_0^\natural=\operatorname{Ker}(\mathrm{id}_{x_0},p)$) soit l'identité, Alors on a $\tau=p$.

\uline{Démonstration} \inserted{de 4.15.7.4.} On est ramené aussitôt, par le \unsure{procédé} \unsure{standart|standard}, au cas où $C=\mathrm{Ens}$, où la vérification est immédiate : les endom. $\tau$ de l'ens $\deleted{\unsure{x|F}}E$ \inserted{$(=x)$} qui satisfont $(a)$ \deleted{sont \unsure{encore|aussi} de la forme} \inserted{correspondent \unsure{biunivoqmt|biunivoquemt}} aux \deleted{\unsure{de|le}} endom. de $E^\natural=pE$, via

\[\left\{\begin{aligned}u&\longmapsto\sigma u\rho\\ \tau&\mathrel{\unsure{\mapstochar\relbar\mapstochar|\longmapsto}}\tau\mid E^\natural\end{aligned}\right.\]

(où \diagram{\begin{tikzcd}[ampersand replacement=\&] E^\natural \arrow[r,hook,shift left,"\sigma"] \& E \arrow[l,shift left,"\rho"] \end{tikzcd}} \unsure{sont les flèches can.|est la flèche can.})

donc $\tau\mid E^\natural=\mathrm{id}\ \Longrightarrow\ u=\sigma\rho=p$.

Choisissons un $j_0\in J$, et considérons les hom.

\[\tau_{j,j_0}^\natural:x_0^\natural\xrightarrow{\sim}x_0^\natural\]

induits par les $\tau_{j,j_0}$ avec $j\geq j_0$. Ici encore, comme l'ens des hom. \deleted{\unsure{poss. en question|pour lesquels}} est de
\scanpage{535}%
\mspage{446}%
% archive sheet number: 534
cardinal $\leq\pi_0$, et $\deleted{\unsure{I|J}}J_{\deleted{\unsure{\geq}}{\geq j_0}}$ est grd devant $\pi_0$, on voit qu'il existe $J'\subset{}_{j_0}\backslash J_{\deleted{\unsure{\geq j_0}}}$, cofinal ds $\mathord{}_{j_0}\backslash J_{\deleted{\unsure{\geq j_0}}}$\inserted{donc ds $J$}, tel que pour $j\in J'$, les $\tau_{j,j_0}^\natural$ soient égaux à un $\deleted{\unsure{=|\tau}}\ \tau':x_0^\natural\xrightarrow{\sim}x_0^\natural$.

Donc si $\deleted{j'\leq j''}\ j',j''\in J'$, avec $j'\leq j''$, on a

\[\underbrace{\tau_{j''j_0}}_{\tau}=\tau_{j''j'}\underbrace{\tau_{j',j}}_{\tau},\]

\[\text{d'où}\qquad\tau_{j''j'}=\mathrm{id}_{x_0^\natural},\]

puisque la flèche $\tau$ \deleted{\unsure{sont des}} est un iso. Dans \deleted{\unsure{si|pour}} le système inductif réindexé par $J'$, \deleted{\unsure{on a}} la condition $4^\circ)$ est \uline{remplacée par la con}dition plus stricte

\[\boxed{4^\circ\text{bis})\quad\deleted{\text{\unsure{pour|que}}}\quad\tau_{j',j}^\natural=\mathrm{id}_{x_0^\natural}.}\]

Mais la condition $\deleted{3}2^\circ)$ devient alors, grâce à 4.15.7.4\inserted{\unsure{et compte tenu de|et condition de} $3^\circ,4$},

\[\boxed{\deleted{3}2^\circ\text{bis})\quad\forall j,\ \text{existe }j'\geq j\ \text{avec }\tau_{j',j}=p}. \]

\deleted{Donc}, OPS $1^\circ,2^\circ\mathrm{bis},3^\circ,4^\circ\mathrm{bis}$. (on désigne à nouveau par $J$, \inserted{\unsure{ledit|ci-dt} $J'$,} l'ens. ordonné grd devant $\pi_0$ indexant.).

\deleted{\unsure{Distinguons deux cas :} \deleted{\unsure{Soit $K$ l'ens ord.}} $a)$ \deleted{$J$} $I$ admet un objet final -- alors $I$ aussi, \unsure{Notons que comme}}

Soit $K$ l'ens ordonné défini par
\scanpage{536}%
\mspage{447}%
% archive sheet number: 535
\[\left|\begin{aligned}&\operatorname{Ob}K=\operatorname{Ob}J\\ &\text{si }j,j'\in\operatorname{Ob}J,\ \text{on a } j\overset{K}{\leq}j'\quad\text{sss}\quad j\overset{J}{=}j',\ \text{ou }j\overset{J}{<}j'\ \text{et }\tau_{j',j}=p\end{aligned}\right.\]

Il est immédiat que c'est une \uline{relation d'ordre} sur $J$\inserted{(car $p^2=p$)}, plus fine que la relation d'ordre initiale.

\deleted{De plus, \deleted{\unsure{on|on a}} : \unsure{Car|car} si $i\overset{K}{<}j$ \deleted{\unsure{i.e.}} (i.e. $i<j$ et $\tau_{ji}=p$) alors}

\deleted{\[\left|\begin{aligned}j\leq j'&\Longrightarrow i\overset{K}{<}j'\\ i'\leq i&\Longrightarrow i'\overset{K}{<}j.\end{aligned}\right.\]}

\deleted{En effet,}

\deleted{\diagram{\begin{tikzcd}[ampersand replacement=\&] x_{\unsure{i|j}} \arrow[r] \arrow[d,equal] \& x_j\arrow[r]\arrow[d,equal] \& x_{j'}\\ x_0 \& \unsure{x|x_0} \& \end{tikzcd}}}

\deleted{on a}

\deleted{\[\begin{aligned}\tau_{j'i}&=\tau_{j'j}\tau_{ji}={}&\tau_{j'j}p\\ &&\Vert\ \text{par }1^\circ\\ &&p\tau_{j'j}\end{aligned}\]}

\deleted{\deleted{\unsure{On en a|On en a a}} donc $p\tau_{j'i}=\tau_{j'j}\unsure{p}$}

Le syst. inductif induit sur $K$ satisfait aux conditions $1^\circ,2^\circ,3^\circ$ de 4.15.7.3. \deleted{Il reste à} D'autre part $K\mathrel{\unsure{\longrightarrow|\relbar}}J$ est cofinal, donc aussi le composé $K\longrightarrow J\longrightarrow I$. Il reste à prouver que $K$ est grde devant $\pi_0$, donc que tte partie $K'\subset K$ telle que $\operatorname{card}K'\leq\pi_0$ admet un majorant dans $K$. Elle admet un majorant \uline{dans $\deleted{\unsure{I|J}}J$}, soit $j$. Par hypothèse 2 bis, il existe $j'\geq j$ tel que $\tau_{j'j}=p$ \deleted{\unsure{donc|et alors}}, et OPS $j'>j$. \deleted{\unsure{Sinon}} (En effet, si $\tau_{j'j}=p$, on aurait \deleted{que}
\scanpage{537}%
\mspage{448}%
% archive sheet number: 536
la forme initiale $2^\circ)$ de la condition 2bis que l'on a $\tau_{j''j}=p$ pour tt $j''\geq j'$, donc si $j$ n'est pas objet final de $J$, OPS $j'>j$. Mais $J$ n'a pas d'objet final, car sinon $I$ en aurait un, cas que l'on pouvait éliminer.) Donc il existe $j'$ tel que $j'\overset K>j$. Or on a ceci (qu'on appliquera au cas $i\in K'$)

\[i\leq j\quad\text{et}\quad j\overset K<j'\ \Longrightarrow\ i\overset K<j'.\]

\diagram{\begin{tikzcd}[ampersand replacement=\&] x_i \arrow[d,equal]\arrow[r,shift left,"{\tau_{j,i}\deleted{\unsure{*}}}"] \& x_j\arrow[l,shift left,"{\unsure{\leftarrow|\deleted{\leftarrow}}}" description,draw=none]\arrow[r,"{\deleted{\unsure{\tau_{j'j}}}\inserted{\tau_{j',j}=p}}"] \& x_{j'} \\ x_0 \& \& \end{tikzcd}}

\[\begin{aligned}\tau_{j'i}&=\tau_{j'j}\tau_{ji}={}&p\tau_{ji}&\overset{?}{=}p\\ &&\underset{\tau_{ji}p}{\Vert}\ \text{par }1^\circ\end{aligned}\]

\[\text{on a en effet}\quad\left.\begin{aligned}\tau_{j',i}&=p\,\tau_{j'i}\\ \tau^{\unsure{\prime|{}}}_{j',i}&=\tau_{j'i}p\end{aligned}\right|\Longrightarrow\underset{\text{par }4.15.7.4}{\tau_{j'i}=p}.\]

Cela achève la démonstration du corollaire 4.15.7.3

\inserted{\uline{Cor.} 4.15.7.5 \inserted{(de 4.15.7.3)} Soit $(x_i)_{i\in I}$ un syst. inductif \inserted{ess. constant,} $I$ grd devant $\pi_0\geq\operatorname{card} C$, dans $C$ \uline{sans projecteurs}. Alors il existe $J\to I$ cofinal, avec $J$ ordonné grd devant $\pi_0$, et tel que le syst. inductif induit sur $J$ soit \uline{constant}.}

\uline{Démontrons le lemme} 4.15.7.1. Il faut expliciter l'existence de deux flèches dans $\mathrm{Ind}(C)$

\diagram{\begin{tikzcd}[ampersand replacement=\&] z \arrow[r,shift left,"\alpha","{\unsure{\deleted{\leftarrow}|{}}}"{description,pos=0.15}] \& \mathop{\text{\textquotedblleft}\lim\text{\textquotedblright}}\limits_{i\in I}x_i=\underline{x}\arrow[l,shift left,"\beta"] \end{tikzcd}}

avec $\beta\alpha=\mathrm{id}_z$, $\alpha\beta=\mathrm{id}_{\underline{x}}$. \deleted{\unsure{La}} La donnée
\scanpage{538}%
\mspage{449}%
% archive sheet number: 537
de $\alpha$ équivaut à celle d'un élément de $\mathop{\unsure{\underline{\lim}|\varinjlim}}\limits_{i\in I}\operatorname{Hom}(z,x_i)$, \deleted{\unsure{cet}} est donc donné par un

\[\alpha_{i_0}:z\longrightarrow x_{i_0},\]

la donnée de $\beta$ équivaut à celle d'un élément de $\varprojlim_{i\in I}\operatorname{Hom}(x_i,z)$, i.e. de flèches

\[\beta_i:x_i\longrightarrow z,\]

avec la condition de compatibilité habituelle. \inserted{Quitte à remplacer $I$ par $I'$, avec $I'\to I$ cofinal,} OPS $I$ ordonné\inserted{(et grd devant $\pi_0$, si $I$ initial l'était)}, et la condition \deleted{\unsure{de}} de compatibilité sur les $\beta_i$ s'écrit

\[\beta_i\tau_{ji}=\beta_j\qquad\text{si }i\leq j\text{ ds }I,\]

où $\tau_{ji}:x_i\mathrel{\unsure{\relbar|\to}}x_j$ désigne la flèche de transition. \deleted{\unsure{D|Il}}Il faut expliciter, en termes de \unsure{nos|ces|ses} données, les conditions

\[\alpha\beta=\mathrm{id}_z,\qquad\beta\alpha=\mathrm{id}_{\deleted{\unsure{z|x}}\underline{x}}\unsure{,|!}\]

\diagram{\begin{tikzcd}[ampersand replacement=\&] z\arrow[r,"{\alpha_{i_0}}"]\arrow[dr] \& x_{i_0}\arrow[r,dash,"{\tau_{i,i_0}}"]\arrow[d,dash,"{\beta_{i_0}}"] \& x_i\arrow[dl,dash,"{\beta_i}"']\\ \& z \& \end{tikzcd}}

La relation $\alpha\beta=\mathrm{id}_z$ équivaut à \deleted{\unsure{$\beta$}}

\begingroup\makeatletter\tagsleft@true\makeatother
\[\deleted{\unsure{\beta z|\alpha\beta}}\ \beta_{i_0}\alpha_{i_0}=\mathrm{id}_z,\tag*{$(*_{i_0})$}\]
\endgroup

elle implique, d'ailleurs, en posant pour $i\geq i_0$

\[\alpha_i\overset{\mathrm{déf}}{=}\tau_{i,i_0}\alpha_{i_0}:z\longrightarrow x_i,\]

que l'on a pour $i\geq i_0$

\begingroup\makeatletter\tagsleft@true\makeatother
\[\beta_i\alpha_i=\mathrm{id}_z\qquad\forall i\geq i_0.\tag*{$(*_i)$}\]
\endgroup
\scanpage{539}
\mspage{450}
% archive sheet number: 538

Posons alors

\[p_i=\alpha_i\beta_i:\unsure{x_i|z_i}\mathrel{\unsure{\relbar|\to}}x_i,\]

il est clair qu'on a

\[\left\{\begin{aligned}&p_i^2=p_i\quad(p_i\text{ un projecteur})\\ &\alpha_i:z\longrightarrow x_i\quad\text{induit un isom.}\\ &\hspace{2em}z\xrightarrow{\sim}x_i^\natural\overset{\mathrm{déf.}}{=}\operatorname{Im}p_i=\operatorname{Ker}(\mathrm{id}_{x_{\unsure{j|i}}},p_i).\end{aligned}\right.\]

L'ens \inserted{ordonné} cofinal cherché $J$ est $i_0\backslash I_{\deleted{\unsure{\geq i_0}}}$, il est filtrant, et est \unsure{grd|gd|pt|petit} devant $\pi_0$ \unsure{si} $I$ l'est. \deleted{Il faut vérifier les conditions $1^\circ$ \unsure{;|à} $4^\circ$.}

À vrai dire, la réduction au cas $I$ ordonné est inutile, on peut procéder directement $J=i_0\backslash I$, Il est immédiat que $\deleted{\unsure{I}}J$ est filtrante ($I$ l'étant) et $i_0\backslash I\mathrel{\relbar}I$ cofinal, enfin que si $I$ est grd devant $\pi_0$\inserted{\unsure{un cardinal infini|un ordre infini}}, il en est de $=$ de $i_0\backslash I$. \deleted{\unsure{Condition|Considérons} \unsure{$1^\circ$|$1^\circ$ :}}

Exprimons que $\beta\alpha=\mathrm{id}_{\underline{x}}$ ; or on voit que $\beta\alpha$ est défini par l'endomorphisme $(p_j)_{j\in J}$ du système inductif $(x_j)_{j\in J}$ --- le fait que ce soit bien un endom. de ce système, i.e. qu'on ait la relation $1^\circ)$ de 4.15.7.1.

\[p_{j'}x_\lambda=x_\lambda p_{\unsure{\lambda|j}}\qquad\text{si }\lambda:j\mathrel{\relbar}j'\text{ ds }J,\]

est immédiat, ça s'écrit

\[\alpha_{j'}\underbrace{\beta_{j'}x_\lambda}_{\beta_j}=\underbrace{x_\lambda\alpha_j}_{\alpha_{j'}}\beta_j\qquad\text{or les deux mbs sont égaux à }\alpha_{j'}\beta_j.\]
\scanpage{540}
\mspage{451}
% archive sheet number: 539

Dire que c'est l'identité, \deleted{\unsure{signifie que}} \inserted{\unsure{est|c'est}} exprimé par la condition $2^\circ)$ sur $(x_j,p_j)_{j\in J}$.

La condition $3^\circ$ est claire par construction, $p_j$ admettant une image isom. à $z$ par $\alpha_j$. D'autre part, on voit de suite que pour une flèche $\lambda:\inserted{j\mathrel{\relbar}j'}$ \unsure{ds|\deleted{de}} $J$, on a commutativité du

\diagram{\begin{tikzcd}[ampersand replacement=\&] x_j\arrow[r,dash,"{x_\lambda}"] \& x_{j'} \\ x_j^\natural\arrow[u,arrows={Hooks[left]-}]\arrow[r,dash,"{x_\lambda^\natural}"] \& x_{j'}^\natural\arrow[u,arrows={Hooks[left]->}] \\ z\arrow[u,"{\overline\alpha_j}"',"\sim"{rotate=90}]\arrow[r,dash,"{\underset{\sim}{\mathrm{id}_z}}"] \& z\arrow[u,"{\overline\alpha_{j'}}"',"\sim"{rotate=90}] \end{tikzcd}},

où $\overline\alpha_j,\overline\alpha_{j'}$ sont les iso. définis par $\alpha_j,\alpha_{j'}$ respectivement. Cela prouve que $x_\lambda^\natural$ est un iso.

Cela prouve que si $(x_i)_{i\in I}$ est ess. constant, il existe $J\mathrel{\relbar} I$ comme indiqué. L'inverse est également vrai, car OPS $J=I$, \unsure{puisque|\deleted{puis}que} $(x_i)_{i\in I}$ et $(x_j)_{j\in J}$ définissent des Ind-objets isomorphes. Alors le système inductif des $x_j^\natural$ a des morphismes de transition des iso, et comme $J$ est \inserted{filtrante donc} 1-connexe, on en déduit \unsure{facilemt|facilement} qu'il est isomorphe
\scanpage{541}%
\mspage{452}%
% archive sheet number: 540
est isomorphe \unsure{:|à} un syst. inductif constant, de valeur $z$. En choisissant un tel isom

\[z\mathrel{\overset{\overline\alpha_j}{\overset{\sim}{\relbar}}}x_{\unsure{j|i}}^\natural,\]

on trouve $z\xrightarrow{\alpha}\underline{x}$, et le système des $p_i$ définit des

\diagram{\begin{tikzcd}[ampersand replacement=\&] x_{\unsure{j|i}}\arrow[r,"{\overline p_i}"]\arrow[rr,bend right=22,dash,"{\beta_i\overset{\mathrm{déf}}{=}\overline\alpha_j^{-1}\,\overline p_i}"'] \& x_i^\natural\arrow[r,"{\overline\alpha_j^{-1}}"] \& \deleted{\unsure{x|z}}z \end{tikzcd}},

et on voit \deleted{\unsure{aussitôt}} \unsure{facilemt|facilement} que l'on aura

\[\beta\alpha=\mathrm{id}_z,\qquad\alpha\beta=\mathrm{id}_{\underline{x}},\qquad\text{qed.}\]

Je \uline{\unsure{peux} maintenant prouver 4.15.3.} C'est en effet la conjonction de 4.15.7.5 (p. 448), et de 4.8.13, qui nous dit \inserted{\unsure{en.|env.}} que si $C$ est \uline{\unsure{Karoubienne|Karoubien}} \unsure{$\nearrow$|*} (condition vérifiée si $C$ est sans projecteurs), alors tt syst. inductif dans $C$, \unsure{grd|petit|pt} devant $\pi_0$, où $\pi_0\geq\operatorname{card}\operatorname{\unsure{red|réd|nerf}}C$, est essentiellement constant.

Qed.
\scanpage{542}
\mspage{453}
% archive sheet number: 541

4.15.8. Je veux revenir ici sur la démonstration inachevée de 4.8.13 (p. \unsure{213|217|211}--216). Je reprends la démonstration où \deleted{\unsure{j|l}} je l'avais laissée (p. 215). Il faut prouver que

\[F(\unsure{y|g|G})=\mathop{\unsure{\underline{\lim}|\varinjlim}}\limits_{i\in I}M_i\]

est de cardinal $\leq\pi_0$. J'ai annoncé qu'on peut se ramener au cas où $M$ est un monoïde de monomorphismes, i.e. que $uv=uv'\Longrightarrow v=v'$ (pour $u,v,v'\in M$). Je vais expliciter cet argument.

Pour tout $i\in I$, soit

\begingroup\makeatletter\tagsleft@true\makeatother
\[N_i=\left\{u\in M\ \middle|\ \begin{aligned}&\text{l'ens }I_{i,u}\text{ des }j\in I\text{ tels que}\\ &j\geq i,\ \alpha_{ji}\deleted{\unsure{\in|=}}=u\quad\text{est cofinal dans }I\\ &\text{i.e. pour tt }j_0\in I,\ \text{existe }j\in I_{i,u}\text{ avec}\\ &j\geq j_0\end{aligned}\right\}.\tag*{(4.15.8.1)}\]
\endgroup

Comme $I$, donc $i\backslash I$ \unsure{grd|grand} devant $\pi_0$, et $\operatorname{card} M\leq\pi_0$, on voit \unsure{aussitôt|immédiatmt} que

\begingroup\makeatletter\tagsleft@true\makeatother
\[N_i\ne\unsure{\varnothing|\phi}.\tag*{(4.15.8.2)}\]
\endgroup

[ Si d'autre part $i'\geq i$, on trouve

\begingroup\makeatletter\tagsleft@true\makeatother
\[N_i=N_{i'}\alpha_{i'i}\tag*{(4.15.8.\unsure{3|8|B})}\]
\endgroup

L'inclusion $N_{i'}\alpha_{i',i}\subset N_i$ est $\pm$ tautologique, \deleted{\unsure{p|f}} prouvons $N_i\subset N_{i'}\alpha_{i'i}$. Soit $u\in N_i$, l'ens \inserted{$J=I_{i,u}\cap i'\backslash I$} des $j\geq i'$ tels que $\alpha_{ji}=u$ est cofinal dans $I$, donc grand devant $\pi_0$.

\diagram{\begin{tikzcd}[ampersand replacement=\&] i\arrow[r]\arrow[rr,bend right=18,"{u=\alpha_{ji}}"',dash] \& i'\arrow[r,"{\alpha_{ji'}}"] \& j\end{tikzcd}}

\deleted{L'application} Donc il existe un $v\in M$ tel que l'ens \inserted{$J'$} des $j\in J$ tels que $\alpha_{ji'}=v$
